\subsection{Invarianza per traslazioni temporali di $\mathcal{H}$}
So che $ S(t_1, t_0) = S(t_1 + \delta, t_0 + \delta) $, nel caso infinitesimo
\[
\mathbbm{1} + i \varepsilon \hat{\mathcal{H}}(t_0) = \mathbbm{1} + i \varepsilon \hat{\mathcal{H}}(t_0 + \delta) \quad \Rightarrow \quad \frac{d}{dt} \hat{\mathcal{H}}(t) = 0
\]
cio\`e $ \hat{\mathcal{H}} \equiv \hat{\mathcal{H}}(\hat{p}, \hat{q}) $ (N.B.: \`e un ``se e solo se'').

Se $ \hat{\mathcal{H}} $ non dipende dal tempo, allora $ S $ e $ \hat{\mathcal{H}} $ commutano. La condizione di partenza equivale a $ S(t_1 + \delta, t_0 + \delta) = S(t_1, t_0) $ con $ \delta $ infinitesimo questa volta,
\begin{align*}
S(t_1 + \delta, t_0 + \delta) &= S(t_1 + \delta, t_1) S(t_1, t_0) S(t_0, t_0 + \delta) = \\
&= S(t_1 + \delta, t_1) S(t_1, t_0) S^{-1}(t_0 + \delta, t_0) = \\
&=  S(t_1 + \delta, t_1) S(t_1, t_0) S^\dag(t_0 + \delta, t_0) = \\
&= \left( \mathbbm{1} + i \delta \hat{\mathcal{H}} \right) S(t_1, t_0) \left( \mathbbm{1} - i \delta \hat{\mathcal{H}} \right) = \\
&= S(t_1, t_0) + i \delta \left[ \hat{\mathcal{H}}, S(t_1, t_0) \right] + \mathcal{O}(\delta) \quad \Rightarrow
\end{align*}
\[
\Rightarrow \quad \left[ \hat{\mathcal{H}}, S(t_1, t_0) \right] = 0.
\]
Inoltre gli autovalori di $ \hat{\mathcal{H}} $ si conservano per traslazione temporale
\[
\hat{\mathcal{H}} \ket{E} = E \ket{E}.
\]
\begin{align*}
0 &= \braket{\varphi|[\hat{\mathcal{H}}, S(t, t_0)]|\psi} = \braket{\varphi|(\hat{\mathcal{H}} S(t, t_0) - S(t, t_0) \hat{\mathcal{H}})|\psi} = \\
&= \int_{-\infty}^{+\infty} \int_{-\infty}^{+\infty} \braket{\varphi|E} \braket{E|(\hat{\mathcal{H}} S - S \hat{\mathcal{H}})|E^\prime} \braket{E^\prime|\psi} \, dE \, dE^\prime = \\
&= \int_{-\infty}^{+\infty} \int_{-\infty}^{+\infty} \braket{\varphi|E} \left( E \braket{E|S|E^\prime} - E^\prime \braket{E|S|E^\prime}\right) \braket{E^\prime|\psi} \, dE \, dE^\prime  = \\
&= \int_{-\infty}^{+\infty} \int_{-\infty}^{+\infty} (E - E^\prime) \braket{\varphi|E} \braket{E|S|E^\prime} \braket{E^\prime|\psi} \, dE \, dE^\prime
\end{align*}
cio\`e $ \braket{E|S|E^\prime} = 0 $ se $ E \neq E^\prime $.

Digressione sulla lagrangiana.
\begin{align*}
\frac{d \mathcal{L}}{dt} \left( \underline{q}(t), \dot{\underline{q}}(t) \right) &= \sum_k \frac{\partial \mathcal{L}}{\partial q_k} \dot{q}_k + \sum_k \frac{\partial \mathcal{L}}{\partial \dot{q}_k} \ddot{q}_k = \left[ \frac{d}{dt} \frac{\partial \mathcal{L}}{\partial \dot{q}_k} = \frac{\partial \mathcal{L}}{\partial q_k} \right] = \\
&= \sum_k \dot{p}_k \dot{q}_k + \sum_k p_k \ddot{q}_k = \frac{d}{dt} \sum_k p_k \dot{q}_k
\end{align*}
\[
\frac{d}{dt} \mathcal{H} \left( \underline{q}(t), \underline{p}(t) \right) = \frac{d}{dt} \left( \sum_k p_k \dot{q}_k - \mathcal{L} \right) = 0.
\]
Fine digressione sulla lagrangiana.

Voglio che $\hat{\mathcal{H}} = -c \mathcal{H}$. Dimensionalmente, poich\'e
\[
S(t, t_0) \ket{\psi(t_0)} = \ket{\psi(t)} \; \xrightarrow[\text{infinitesim.}] \; \ket{\psi(t)} = e^{(t - t_0) \frac{d}{dt}} \ket{\psi(t)} |_{t_0} =
\]
\[
= \left( \mathbbm{1} + (t - t_0) \frac{d}{dt} + \mathcal{O}(t - t_0) \right) \ket{\psi(t_0)} \quad \Rightarrow
\]
\[
\Rightarrow \quad \hat{\mathcal{H}}(t) \ket{\psi(t)} = \frac{d}{dt} \ket{\psi(t)}
\]
la costante \`e un'azione, quindi $ c \equiv \frac{1}{\hbar} $. Riscrivo
\[
\hat{\mathcal{H}} = - \frac{1}{\hbar} \mathcal{H} \left( \hat{p}, \hat{q} \right) \quad \Rightarrow \quad \boxed{i \hbar \frac{d}{dt} \ket{\psi(t)} = \mathcal{H}( t; \hat{p}, \hat{q}) \ket{\psi(t)}} \; \text{eq. di Schr\"odinger.}
\]
Per $ \mathcal{H} = T + V = \frac{\hat{p}^2}{2m} + V(\hat{q}) $ per stati $ \ket{q} $
\[
\braket{q|i \hbar \frac{d}{dt}|\psi(t)} = \braket{q|\frac{\hat{p}^2}{2m} + V(\hat{q})|\psi(t)} \quad \Rightarrow
\]
\[
\left[
\begin{aligned}
\braket{q|\hat{p}^2|\psi} &= \int_{-\infty}^{+\infty} \braket{q|\hat{p}|k} \braket{k|\hat{p}|\psi} \, dk = \\
&= \int_{-\infty}^{+\infty} (\hbar k)^2 \braket{q|k} \braket{k|\psi} \, dk = \\
&= \int_{-\infty}^{+\infty} \frac{(\hbar k)^2}{\sqrt{2\pi}} e^{iqk} \psi(k) \, dk \\
\braket{q|\hat{q}|\psi} &= q \psi(q) \\
\braket{q|F(\hat{q})|\psi} &= F(q) \, \psi(q)
\end{aligned}
\right]
\]
\[
\Rightarrow \quad i \hbar \frac{d \psi}{dt}(q; t) = - \frac{\hbar^2}{2m} \frac{d^2 \psi}{dq^2}(q;t) + V(q) \psi(q;t).
\]

Per l'evoluzione temporale, $ \ket{\psi(t)} = S(t, t_0) \ket{\psi(t_0)} $, quindi
\[
i \hbar \frac{d}{dt} S(t, t_0) \ket{\psi(t_0)} = \mathcal{H}(t) S(t, t_0) \ket{\psi(t_0)}
\]
e, a seconda dei casi, ho diverse discussioni.

Caso 1.
\[
\frac{d\mathcal{H}}{dt} = 0 \qquad \text{(stato legato statico)}
\]
\[
i \hbar \frac{d}{dt} S(t, t_0) = \mathcal{H} S(t, t_0)
\]
\[
\boxed{S(t, t_0) = \exp \left\{ \frac{t - t_0}{i \hbar}\mathcal{H} \right\} }
\]

Caso 2.
\[
\frac{d\mathcal{H}}{dt} \neq 0 \text{ e } [\mathcal{H}(t), \mathcal{H}(t^\prime)] = 0
\]
\[
i \hbar \frac{d}{dt} S(t, t_0) = \mathcal{H}(t) S(t, t_0) \quad \Rightarrow \quad \boxed{S(t, t_0) = \exp{\left\{ \frac{1}{i \hbar} \int_{t_0}^{t} \mathcal{H}(t^\prime) \, dt^\prime \right\}}}
\]

Caso 3.
\[
\frac{d\mathcal{H}}{dt} \neq 0 \text{ e } [\mathcal{H}(t), \mathcal{H}(t^\prime)] \neq 0
\]
\begin{figure}
\centering
\begin{tikzpicture}
  \draw[thick] (0,5) -- (0,0) -- (5,0);
  \draw[dashed] (0,0) -- (5,5);
  \draw (0,3) -- (3,3) -- (3,0);

  \node (t0) at (-.2,-.2) {$t_0$};
  \node (t1) at (3,-.3) {$t'$};
  \node (t2) at (-.3,3) {$t''$};
\end{tikzpicture}
\caption{Logica di funzionamento della funzione $T[O_1(t_1)O_2(t_1)]$.}
\label{schema_T}
\end{figure}
Definisco
\[
T \left[ O_1(t_1) O_2(t_1) \right] \equiv \left\{
\begin{aligned} 
    & O_1(t_1) O_2(t_1) & t_1 > t \\ 
    & O_2(t_1) O_1(t_1) & t_1 < t
\end{aligned} \right. .
\]
Sviluppando al secondo ordine, facendo uso dello schema in Figura \ref{schema_T},
\begin{align*}
S(t, t_0) &= T \exp{\left\{ \frac{1}{i \hbar} \int_{t_0}^t \mathcal{H}(t^\prime) \, dt^\prime \right\}} = \\
&= \mathbbm{1} + \frac{1}{i \hbar} \int_{t_0}^t \mathcal{H}(t^\prime) \, dt^\prime + \frac{1}{2} \frac{1}{(i \hbar)^2} \int_{t_0}^t \int_{t_0}^t T \left( \mathcal{H}(t^\prime) \mathcal{H}(t'') \right) \, dt' \, dt'' = \\
&= \mathbbm{1} + \frac{1}{i \hbar} \int_{t_0}^t \mathcal{H}(t^\prime) \, dt^\prime + \frac{1}{2} \frac{1}{(i \hbar)^2} \left( \int_{t_0}^t \int_{t_0}^{t'} \mathcal{H}(t') \mathcal{H}(t'') \, dt' \, dt'' + \right. \\
& \left. \quad + \int_{t_0}^t \int_{t'}^t \mathcal{H}(t'') \mathcal{H}(t') \, dt'' \, dt' \right) = \\
&= \mathbbm{1} + \frac{1}{i \hbar} \int_{t_0}^t \mathcal{H}(t^\prime) \, dt^\prime + \frac{1}{(i \hbar)^2} \int_{t_0}^t \int_{t_0}^{t'} \mathcal{H}(t') \mathcal{H}(t'') \, dt' \, dt''.
\end{align*}
Quindi per l'$n$-esimo termine dello sviluppo in serie di Taylor dell'esponenziale applico lo stesso metodo, cio\`e
\[
T \left( \int_{t_0}^t \mathcal{H}(t_1) \, dt_1 \right) \left( \int_{t_0}^{t_1} \mathcal{H}(t_2) \, dt_2 \right) \cdots \left( \int_{t_0}^{t_{n - 1}} \mathcal{H}(t_n) \, dt_n \right) =
\]
\[
= n! \int_{t_0}^t \int_{t_0}^{t_1} \int_{t_0}^{t_2} \cdots \int_{t_0}^{t_{n - 1}} \mathcal{H}(t_1) \mathcal{H}(t_2) \cdots \mathcal{H}(t_n) \, dt_1 \, dt_2 \cdots dt_n
\]
e gli $ n! $ si semplificano. Adesso controllo se $ S(t, t_0) $ \`e soluzione dell'equazione di Schr\"odinger.
\begin{align*}
i \hbar \frac{d}{dt} S(t, t_0) &= \mathcal{H}(t) + \frac{1}{i \hbar} \mathcal{H}(t) \int_{t_0}^t \mathcal{H}(t_2) \, dt_2 + \cdots + \\
&\quad + \frac{1}{(i \hbar)^{n + 1}} \mathcal{H}(t) \int_{t_0}^t \int_{t_0}^{t_2} \cdots \int_{t_0}^{t_{n - 1}} \mathcal{H}(t_2) \cdots \mathcal{H}(t_n) \, dt_2 \cdots dt_n = \\
&= \mathcal{H}(t) S(t, t_0)
\end{align*}

Nel caso pi\`u semplice in cui $ \frac{d}{dt} \mathcal{H} = 0 $ so che
\[
S(t, t_0) = \exp\left\{ \frac{1}{i \hbar} (t - t_0) \mathcal{H} \right\}
\]
allora l'operatore di evoluzione temporale lo posso calcolare solo quando $ \mathcal{H} \ket{n} = E_n \ket{n} $ (diagonalizzabilit\`a dell'operatore), da cui poi $ \braket{n|\psi} = c_n $ e $ p_n = |c_n|^2$.

Se ho un potenziale per cui so calcolare gli stati di energia, posso calcolare gli elementi di matrice, ottenendo
\[
S(t, t_0) = \sum_{n,m} \ket{m} \braket{m|S(t, t_0)|n} \bra{n} =
\]
\[
\left[ \exp{\left\{ \frac{1}{i \hbar} (t - t_0)  \mathcal{H} \right\} } \ket{n} = \exp{\left\{ \frac{1}{i \hbar} (t - t_0) E_n \right\}} \ket{n} \right]
\]
\[
= \sum_{n,m} \ket{m} \braket{m|\exp{\left\{ \frac{1}{i \hbar} (t - t_0) E_n \right\}}|n} \bra{n} =
\]
\[
= \sum_{n,m} \exp{\left\{ \frac{1}{i \hbar} (t - t_0) E_n \right\}} \ket{m} \braket{m|n} \bra{n} = 
\]
\[
= \sum_n \exp{\left\{ \frac{1}{i \hbar} (t - t_0) E_n \right\}} \ket{n} \bra{n}
\]
operatore diagonale nella base degli $ \ket{n} $.

Quindi per $ \ket{\psi(t_0)} = \ket{k} $ al tempo $ t $
\[
\ket{\psi(t)} = S(t, t_0) \ket{k} = \sum_n \exp{\left\{ \frac{1}{i \hbar} (t - t_0) E_n \right\}} \ket{n} \braket{n|k} = e^{\frac{1}{i \hbar}(t-t_0)E_k}\ket{k}
\]
cio\`e $ \ket{\psi(t_0)} $ e $ \ket{\psi(t)} $ \textbf{differiscono per una fase}. Inoltre
\[
\ket{\psi(t_0)} = \sum_k \ket{k} \overbrace{\braket{k|\psi(t_0)}}^{c_k(t_0)} = \sum_k c_k(t_0) \ket{k} \quad \Rightarrow
\]
\[
\Rightarrow \quad \ket{\psi(t)} = \sum_n \exp{\left\{ \frac{1}{i \hbar} (t - t_0) E_n \right\}} \ket{n} \overbrace{\braket{n|\psi(t)}}^{c_n(t)} \; \Rightarrow
\]
\[
\Rightarrow \; c_n(t) = c_n(t_0) \exp{\left\{ \frac{1}{i \hbar} (t - t_0) E_n \right\}}
\]
cio\`e la probabilit\`a si conserva perch\'e 
\[
p_n(t) = |c_n(t)|^2 = |c_n(t_0)|^2 = p_n(t_0).
\]

Per lo stato $ \ket{\psi} $ il valor medio di un operatore sulla base degli $ \ket{n} $ autostati dell'energia 
\[
\braket{n|A|n} = \braket{n(t_0)|e^{- \frac{1}{i \hbar} (t_n - t_0) E_n} A e^{\frac{1}{i \hbar} (t_n - t_0) E_n}|n(t_0)} = \braket{n(t_0)|A|n(t_0)}.
\]
Gli stati $ \ket{n} $ sono detti \textbf{stati stazionari}.

\subsubsection{Esempio: hamiltoniana e matrici di Pauli}
$$ \mathcal{H} = B \sigma_z, \qquad \sigma_z = \begin{pmatrix}
1 & 0 \\
0 & -1
\end{pmatrix}.
\[
\]
S = \exp \left\{ t \frac{B \sigma_z}{i\hbar} \right\}
\[
\]
\left| \braket{\psi_1|\psi_2} \right|^2 = \left| \braket{\psi_1|\psi_2(t)} \right|^2 = \left| \braket{\psi_1|S|\psi_2(0)} \right|^2
\[
\]
\ket{\psi(t)} = \exp \left\{ \frac{Bt\sigma_z}{i\hbar} \right\} \frac{\ket{+} + \ket{-}}{\sqrt{2}} = \frac{1}{\sqrt{2}} \left( e^{\frac{Bt}{i\hbar}} \ket{+} + e^{-\frac{Bt}{i\hbar}} \ket{-} \right)
\[
\]
\bra{\psi_2} = \frac{1}{\sqrt{2}} \left( \bra{+} + \bra{-} \right)
\[
\]
\left| \braket{\psi_2|\psi_1} \right|^2 = \frac{1}{4} \left| e^{\frac{Bt}{i\hbar}} + e^{-\frac{Bt}{i\hbar}} \right|^2 = \cos^2 \frac{Bt}{\hbar}
\[

\subsubsection{Esempio: ancora su hamiltoniana e matrici di Pauli}
\]
\mathcal{H} = t B \sigma_z.
\[
\]
S = \exp \left\{ \frac{1}{i\hbar} \int_{0}^t \mathcal{H}(t') \, dt' \right\} = \exp \left\{ \frac{t^2 B \sigma_z}{2i\hbar} \right\}.
\[
\]
\ket{\psi(t)} = \exp \left\{ \frac{Bt^2\sigma_z}{i\hbar} \right\} \frac{\ket{+} + \ket{-}}{\sqrt{2}} = \frac{1}{\sqrt{2}} \left( e^{\frac{Bt^2}{i\hbar}} \ket{+} + e^{-\frac{Bt^2}{i\hbar}} \ket{-} \right)
\[
\]
\left| \braket{\psi_2|\psi_1} \right|^2 = \frac{1}{4} \left| e^{\frac{Bt^2}{i\hbar}} + e^{-\frac{Bt^2}{i\hbar}} \right|^2 = \cos^2 \frac{Bt^2}{\hbar}
$$